\documentclass[11pt]{article}

\usepackage[margin=1in]{geometry}
\usepackage{amsmath}
\usepackage{booktabs}
\usepackage{array}
\usepackage{hyperref}
\usepackage{xcolor}
\usepackage{enumitem}
\usepackage{caption}

\hypersetup{
    colorlinks=true,
    linkcolor=blue!50!black,
    urlcolor=blue!50!black,
    citecolor=blue!50!black
}

\title{A FHIR-to-OMOP Common Data Model ETL Pipeline: \\
Design, Verification, and a Decomposed Concept-Mapping \\
Coverage Metric on the MIMIC-IV-FHIR Demo Dataset}

\author{Saileshwar Dumpala \\ \small \texttt{dsailesh2001@gmail.com}}

\date{September 2026}

\begin{document}

\maketitle

\begin{abstract}
Electronic health record systems increasingly expose clinical data through
Fast Healthcare Interoperability Resources (FHIR), an interoperability
standard designed for point-of-care data exchange, while large-scale
observational health research relies on the Observational Medical Outcomes
Partnership (OMOP) Common Data Model (CDM), a research-oriented schema
designed for standardized, cross-site analysis. This report describes and
verifies an open, end-to-end extract-transform-load (ETL) pipeline that maps
FHIR clinical resources from the MIMIC-IV Clinical Database Demo on FHIR
(100 de-identified patients) into OMOP CDM v5.4, covering the
\texttt{PERSON}, \texttt{VISIT\_OCCURRENCE}, \texttt{CONDITION\_OCCURRENCE},
and \texttt{MEASUREMENT} tables, with terminology mapping through the OHDSI
standardized vocabularies. I report three contributions beyond the
implementation itself. First, I show that a single blended
concept-mapping-coverage statistic can misrepresent pipeline correctness,
and I present a decomposed metric that separates source-data coverage
ceilings from mapping-logic failures, distinguishing a case where low
coverage (1.11\%) reflects a true source-data limitation, fully verified
against two independent data releases, from a case (96.5\%) where nearly all
records carry a mappable code. Second, I describe a synthetic-fixture
testing methodology developed specifically to satisfy continuous-integration
(CI) testing under a real constraint common to OMOP-based projects: OHDSI
Athena vocabulary files require individual license acceptance and cannot be
redistributed, so CI cannot use real vocabulary content. Third, I document
two data-integrity edge cases surfaced by deliberately adversarial fixture
design -- absent from the real 637-encounter dataset by chance rather than
by design -- one of which was fixed at the SQL level and one of which is a
structural limitation documented with a regression-guarding test rather than
silently patched. All quantitative claims in this report were verified
against a running instance of the pipeline, not assumed from code review.
\end{abstract}

\section{Introduction}

Fast Healthcare Interoperability Resources (FHIR) has become the dominant
standard for exchanging clinical data between systems at the point of
care~\cite{hl7fhir}. Observational health research, by contrast, is
increasingly organized around the OMOP Common Data Model, maintained by the
Observational Health Data Sciences and Informatics (OHDSI) collaborative,
which standardizes both schema and terminology across participating sites
so that analytic code can run unmodified against data collected under
different local systems~\cite{hripcsak2015ohdsi}. These two standards serve
different purposes -- interoperability versus standardized analysis -- and
an increasing amount of practical data engineering work in health
informatics consists of translating from the former into the latter.

This report documents an implementation of that translation for a bounded,
verifiable scope: four FHIR resource types (\texttt{Patient},
\texttt{Encounter}, \texttt{Condition}, \texttt{Observation}) mapped into
four corresponding OMOP CDM v5.4 tables (\texttt{PERSON},
\texttt{VISIT\_OCCURRENCE}, \texttt{CONDITION\_OCCURRENCE},
\texttt{MEASUREMENT}), built and verified against the MIMIC-IV Clinical
Database Demo on FHIR, a 100-patient, publicly available de-identified
derivative of MIMIC-IV~\cite{bennett2023mimicfhir, bennett2025physionet,
johnson2023mimiciv}. I treat this as an engineering and methodology report
rather than a claim of novel algorithmic contribution: the value of this
work is in a fully open, reproducible implementation, an honest accounting
of where the pipeline's outputs are and are not trustworthy, and two
specific methodological points -- coverage-metric decomposition and
license-constrained CI design -- that generalize to other OMOP ETL efforts
built on top of licensed OHDSI vocabularies.

The complete source code, verification queries, and a continuous-integration
suite are publicly available; see Section~\ref{sec:availability}.

\section{Background}

\subsection{FHIR and OMOP}

FHIR represents clinical information as discrete, linked \emph{resources}
(e.g., \texttt{Patient}, \texttt{Encounter}, \texttt{Condition},
\texttt{Observation}), each typically serialized as JSON and referencing
related resources by identifier. FHIR's design optimizes for point-of-care
exchange: a single patient's record, retrieved and updated as care happens.

OMOP CDM, in contrast, is a fixed relational schema (39 tables in v5.4)
designed so that identical analytic code -- cohort definitions, data
characterization, statistical models -- can run against any conformant
database regardless of the source EHR system. Every clinically meaningful
column in OMOP is backed by a \emph{concept}: a row in the \texttt{CONCEPT}
table drawn from a standardized vocabulary (SNOMED CT for conditions, LOINC
for laboratory and clinical measurements, and others), rather than a free-
text or locally-coded value. Resolving a source system's local codes to
OMOP standard concepts -- \emph{concept mapping} -- is therefore a
required, non-optional step of any FHIR-to-OMOP pipeline, not an optional
enrichment.

\subsection{Standard vs.\ source concepts and \texttt{CONCEPT\_RELATIONSHIP}}

OMOP distinguishes a \emph{source concept} (what the original data said --
an ICD-9-CM or ICD-10-CM code, a LOINC code, a locally defined code) from a
\emph{standard concept} (OMOP's canonical concept for that clinical fact).
Every relevant OMOP table stores both (e.g.\ \texttt{condition\_concept\_id}
alongside \texttt{condition\_source\_concept\_id}), so no information is
lost in translation. The mechanism that resolves a source concept to a
standard concept is a lookup in \texttt{CONCEPT\_RELATIONSHIP} where
\texttt{relationship\_id = `Maps to'} -- not a formula, but a large,
OHDSI-curated crosswalk table distributed as part of the Athena vocabulary
release.

\subsection{Dataset}

The MIMIC-IV Clinical Database Demo on FHIR (v2.1.0) is a 100-patient,
publicly accessible, de-identified subset of MIMIC-IV re-expressed as FHIR
R4 resources~\cite{bennett2023mimicfhir, bennett2025physionet}, itself
derived from the base MIMIC-IV Clinical Database
Demo~\cite{johnson2023mimiciv}, both distributed via
PhysioNet~\cite{goldberger2000physionet}. This pipeline consumes the FHIR
release directly; the base (non-FHIR) demo was used only as an independent
cross-check for one data-coverage claim (Section~\ref{sec:coverage}).

\section{System Design}

\subsection{Schema separation}

The pipeline uses two PostgreSQL schemas: \texttt{staging}, which mirrors
the source FHIR resource shape (one JSONB column per resource type, raw and
disposable), and \texttt{cdm}, the standard OMOP v5.4 relational schema.
Keeping these separate means the staging layer can be dropped and reloaded
without affecting anything downstream, and no analytic tool is ever pointed
at unvalidated, source-shaped data.

\subsection{Identity resolution}

FHIR resource identifiers are string UUIDs; every OMOP primary key is an
integer surrogate. A small set of crosswalk tables (e.g.\
\texttt{staging.person\_id\_map}, \texttt{staging.visit\_id\_map}), each a
\texttt{GENERATED ALWAYS AS IDENTITY} table populated with
\texttt{ON CONFLICT DO NOTHING} semantics, assigns a stable integer to each
UUID exactly once. This makes every load in the pipeline idempotent: re-
running any stage against a database that already contains data neither
duplicates records nor reassigns identifiers.

\subsection{Scope}

Four FHIR resources map to four OMOP tables (Table~\ref{tab:scope}).
\texttt{Location}, \texttt{Organization}, and \texttt{Provider} were
explicitly out of scope; \texttt{location\_id}, \texttt{provider\_id}, and
\texttt{care\_site\_id} remain \texttt{NULL} throughout. Medication
resources (\texttt{DRUG\_EXPOSURE}) were deferred as future work
(Section~\ref{sec:future}).

\begin{table}[h]
\centering
\caption{Pipeline scope: FHIR resource to OMOP CDM table.}
\label{tab:scope}
\begin{tabular}{ll}
\toprule
FHIR Resource & OMOP CDM Table \\
\midrule
Patient & \texttt{PERSON} \\
Encounter & \texttt{VISIT\_OCCURRENCE} \\
Condition & \texttt{CONDITION\_OCCURRENCE} \\
Observation & \texttt{MEASUREMENT} \\
\bottomrule
\end{tabular}
\end{table}

\section{Implementation}

\subsection{PERSON and VISIT\_OCCURRENCE}

\texttt{Patient} resources map to \texttt{cdm.person}, resolving gender,
race, and ethnicity against a small, fixed set of OMOP reference concepts
(e.g.\ concept ID 8507 for male, 8527 for White, 38003564 for Not Hispanic
or Latino); where no US Core race/ethnicity extension is present, or a
given code (e.g.\ \texttt{ASKU}, ``asked but unknown'') has no OMOP
equivalent, the corresponding concept ID is left at \texttt{0} rather than
guessed. \texttt{Encounter} resources map to \texttt{cdm.visit\_occurrence};
encounters missing \texttt{period.end} are excluded, since
\texttt{visit\_end\_date} is \texttt{NOT NULL} in the OMOP DDL and no valid
value can be synthesized.

\subsection{CONDITION\_OCCURRENCE and MEASUREMENT}

\texttt{Condition} resources have no date field of their own in this
dataset; \texttt{condition\_start\_date} is instead borrowed from the
linked visit's \texttt{visit\_start\_date}, a judgment call made explicit
in the pipeline's documentation rather than left implicit in the SQL.
ICD-9-CM/ICD-10-CM source codes are reformatted to match
\texttt{cdm.concept.concept\_code} conventions (a decimal point inserted
after the third character, e.g.\ \texttt{5715} $\to$ \texttt{571.5}) and
resolved to standard SNOMED concepts via \texttt{CONCEPT\_RELATIONSHIP}.

\texttt{Observation} resources map to \texttt{cdm.measurement}.
\texttt{measurement\_date} is taken directly from \texttt{effectiveDateTime}
(unlike \texttt{Condition}, \texttt{Observation} carries its own date).
Value handling covers three shapes present in the data:
\texttt{valueQuantity} (numeric, with optional \texttt{referenceRange}),
\texttt{valueString}, and \texttt{valueCodeableConcept} display text as a
fallback. \texttt{value\_source\_value} is truncated to 50 characters via
\texttt{LEFT(...,\ 50)} to satisfy the corresponding \texttt{varchar(50)}
schema constraint -- a deliberate, documented lossy step rather than a
silent one; the untruncated text is not preserved elsewhere in the
pipeline. Observations are resolved to standard concepts only when tagged
with a coding whose \texttt{system} is \texttt{http://loinc.org}; other
coding systems present in the source data have no standard-vocabulary
target in this pipeline's scope.

\section{Concept-Mapping Coverage: A Decomposed Metric}
\label{sec:coverage}

A single concept-mapping coverage percentage conflates two structurally
different situations: a source record that simply never carried a
mappable code, and a source record whose code exists but has no
corresponding entry in \texttt{CONCEPT\_RELATIONSHIP}. Reporting only a
blended figure can misrepresent which of these is occurring, with direct
consequences for how a reader should interpret pipeline correctness.

For \texttt{CONDITION\_OCCURRENCE}, this distinction does not materially
matter: staged condition records overwhelmingly carry an ICD-9-CM or
ICD-10-CM code, so a single ratio is already representative. Of 5{,}051
loaded rows, 4{,}872 (96.5\%) resolved to a standard SNOMED concept; the
unmapped 179 rows (3.5\%) are predominantly administrative or status codes
(e.g.\ ``Do not resuscitate status'') with no corresponding disease
concept, not resolution failures.

For \texttt{MEASUREMENT}, the distinction is material. Of 813{,}540 loaded
observations, only 9{,}042 (1.11\%) carry a coding tagged
\texttt{system\ =\ http://loinc.org}; the remaining 98.9\% use MIMIC's own
local chartevents/labevents item dictionaries. This was confirmed against
two independent sources -- the FHIR export's coding systems and the base
(non-FHIR) MIMIC-IV Clinical Database Demo's \texttt{d\_labitems}/
\texttt{d\_items} tables -- neither of which provides a LOINC crosswalk for
these local dictionaries. Reported alone, 1.11\% coverage could plausibly
be read as a mapping-logic defect. Decomposing the statistic shows
otherwise: of the 9{,}042 rows that carried a real LOINC code, all 9{,}042
(100\%) resolved successfully via \texttt{CONCEPT\_RELATIONSHIP}. The
mapped-row count equals the LOINC-coded-row count exactly. The 1.11\%
figure is therefore a source-data coverage ceiling, verified rather than
assumed, and not evidence of a defect in the mapping implementation.

I suggest this decomposition -- coverage among rows carrying a
vocabulary-eligible code, reported separately from overall table coverage
-- as a general practice for OMOP ETL reporting whenever a source
system's code coverage for a given domain and vocabulary is not close to
complete, since a single blended number is otherwise ambiguous between two
causes with very different engineering implications.

\section{Testing Methodology}
\label{sec:testing}

\subsection{The licensing constraint}

OHDSI Athena vocabulary files require individual license acceptance prior
to download and are not licensed for redistribution. This rules out both
of the obvious approaches to CI testing: automating a real vocabulary
download (no scriptable license acceptance, and no right to redistribute
the result even if it could be scripted), and committing any slice of a
real \texttt{CONCEPT}/\texttt{CONCEPT\_RELATIONSHIP} extract to a public
repository.

\subsection{Synthetic fixture design}

The pipeline's CI suite instead constructs a small, entirely synthetic
fixture: fabricated FHIR resources (3 patients, 4 encounters, 4 conditions,
7 observations) and a structurally valid but fabricated OMOP vocabulary
subset, run through the pipeline's real, unmodified code. The vocabulary
fixture distinguishes two categories of concept differently. Fixed
reference concepts that pipeline SQL hardcodes by numeric ID (gender, race,
ethnicity, visit type, type concept, and concept ID 0) use their real OMOP
IDs, since these are stable identifiers referenced directly by the
transform logic, not the copyrighted curated content of a vocabulary.
Source and standard clinical concepts (ICD-9-CM/ICD-10-CM/SNOMED/LOINC) are
fully fabricated, since these are resolved dynamically by code lookup
rather than hardcoded, and fabricating them avoids any licensing concern
entirely.

The fixture was deliberately constructed to include adversarial cases
rather than a happy path alone: one encounter is missing
\texttt{period.end} (excluding it from \texttt{VISIT\_OCCURRENCE} by
design), and one condition and one observation are linked to that
excluded encounter -- a combination never exercised by the real
637-encounter dataset, in which zero encounters were missing a date.

\subsection{Two edge cases found}

Running the adversarial fixture surfaced two distinct behaviors, handled
differently based on the OMOP schema's own constraints.

\textbf{MEASUREMENT (fixed).} \texttt{visit\_occurrence\_id} was originally
resolved from the visit ID crosswalk alone, which assigns an identifier to
every encounter unconditionally, regardless of whether that encounter's
visit actually loaded into \texttt{cdm.visit\_occurrence}. Against the
fixture, this produced a foreign-key violation that aborted the entire
multi-row \texttt{INSERT...SELECT} statement -- meaning a single excluded
encounter could have zeroed out all 813{,}540 real measurement rows had one
existed in the production data. The fix changes the join against
\texttt{cdm.visit\_occurrence} to a \texttt{LEFT JOIN}, selecting
\texttt{NULL} when the visit never loaded. This is valid because
\texttt{MEASUREMENT.visit\_occurrence\_id} is nullable and
\texttt{MEASUREMENT} carries its own date independent of the visit. The fix
was reverified as a no-op against the real Phase 3 load (zero of 637 real
encounters were excluded, so the join always resolved).

\textbf{CONDITION\_OCCURRENCE (documented limitation, not fixed).} The
analogous case for conditions silently drops the row instead of failing:
no error, no log entry, simply absent from \texttt{cdm.condition\_occurrence}.
This cannot be fixed the same way, because \texttt{condition\_start\_date}
is \texttt{NOT NULL} and \texttt{CONDITION\_OCCURRENCE} has no date field
of its own to fall back on -- it borrows the visit's date by design
(Section 4.2). A \texttt{LEFT JOIN} here would only convert a foreign-key
failure into a \texttt{NOT NULL} failure. This is documented as a genuine
structural limitation of the current design, not patched, and guarded by a
regression test asserting the exact expected row count so that a future
change to this behavior is noticed rather than passing silently.

\subsection{CI architecture}

Continuous integration (GitHub Actions) provisions PostgreSQL 16 via a
native \texttt{services:} container, applies the OMOP v5.4 DDL, loads the
fixture through the same scripts and SQL used against real data, and
asserts 20 exact expected values (row counts, per-patient
gender/race/ethnicity resolution, both edge-case behaviors above,
concept-mapping counts, truncation, and referential integrity) via a
dedicated test script. All environment-specific values (database host,
port, credentials, source data directory) are parameterized via standard
\texttt{PG*} environment variables read natively by \texttt{psql}, so local
and CI execution share identical code paths.

\section{Results}

Table~\ref{tab:results} summarizes verified pipeline output on the real
100-patient dataset. Every figure was confirmed by direct query against a
populated database, not derived from code inspection alone.

\begin{table}[h]
\centering
\caption{Verified pipeline output, MIMIC-IV Clinical Database Demo on FHIR.}
\label{tab:results}
\begin{tabular}{lrl}
\toprule
Table & Rows loaded & Notes \\
\midrule
\texttt{cdm.person} & 100 & 0 orphaned references \\
\texttt{cdm.visit\_occurrence} & 637 & 0 orphaned \texttt{person\_id} \\
\texttt{cdm.condition\_occurrence} & 5{,}051 & 96.5\% concept-mapped \\
\texttt{cdm.measurement} & 813{,}540 & 1.11\% mapped; 100\% within \\
& & the LOINC-coded subset \\
\bottomrule
\end{tabular}
\end{table}

CI (GitHub Actions) passes against the synthetic fixture on every push,
confirmed via an actual GitHub Actions run rather than local simulation
alone. One CI-only failure occurred during development: two shell scripts
lost their git-tracked executable bit on a Windows/OneDrive-backed
checkout, causing an exit-126 permission error when invoked as
\texttt{./script.sh}. This was fixed at the git-index level
(\texttt{git update-index --chmod=+x}) and defensively, by invoking scripts
via \texttt{bash script.sh} rather than relying on the executable bit at
all.

\subsection{Demonstration query}

As a concrete illustration that the loaded CDM supports genuine
cross-table analytic questions, I ran a query joining
\texttt{CONDITION\_OCCURRENCE}, \texttt{CONCEPT}, \texttt{VISIT\_OCCURRENCE},
and \texttt{MEASUREMENT} to compute, for each standard condition concept,
the number of distinct patients diagnosed with it and their average count
of measurements recorded during the diagnosing visit. Table~\ref{tab:demo}
shows the top results. The leading chronic-disease conditions (hypertension,
hyperlipidemia, type 2 diabetes) are consistent with expectations for an
ICU-derived cohort, and average measurement intensity per visit is higher
for acuity-associated conditions (acidosis, thrombocytopenia, acute kidney
injury) than for chronic-but-stable conditions (hypertension, tobacco
dependence), a pattern consistent with clinical plausibility though not
independently validated against ground truth.

\begin{table}[h]
\centering
\small
\caption{Demonstration query output (top 8 of 20 by patient count).}
\label{tab:demo}
\begin{tabular}{lrrr}
\toprule
Condition & Patients & Diagnoses & Avg.\ meas./visit \\
\midrule
History of event & 72 & 320 & 301.0 \\
Essential hypertension & 55 & 134 & 186.9 \\
Hyperlipidemia & 47 & 114 & 267.3 \\
Acute kidney injury & 32 & 77 & 440.7 \\
Type 2 diabetes mellitus & 30 & 102 & 224.2 \\
Atrial fibrillation & 27 & 68 & 363.1 \\
Anemia & 26 & 45 & 401.2 \\
Urinary tract infectious disease & 23 & 47 & 318.9 \\
\bottomrule
\end{tabular}
\end{table}

\section{Limitations}

I state these plainly rather than minimize them. \texttt{MEASUREMENT}
concept-mapping coverage (1.11\%) is a hard ceiling imposed by the source
data, not something further mapping logic can improve without a different
data source or a manually constructed local-code-to-LOINC crosswalk, which
was out of scope for this work. A condition whose linked encounter is
missing a start or end date is silently absent from
\texttt{CONDITION\_OCCURRENCE}, a behavior never observed in the real
100-patient dataset (zero such cases existed) but present by construction
and guarded by a regression test. \texttt{LOCATION}, \texttt{ORGANIZATION},
and \texttt{PROVIDER} are not modeled. CI validates pipeline
\emph{code} correctness against known synthetic inputs; it cannot and does
not validate real-world concept-mapping coverage or data quality, which
were verified separately against the real load. \texttt{DRUG\_EXPOSURE} is
not implemented.

\section{Future Work}
\label{sec:future}

Planned extensions include \texttt{DRUG\_EXPOSURE} via RxNorm; a
comorbidity co-occurrence network over \texttt{CONDITION\_OCCURRENCE}
(nodes as standard condition concepts, edges weighted by shared-patient
count), as a first network-analysis layer over the verified CDM data; a
predictive model (e.g.\ readmission or in-hospital mortality risk) using
\texttt{PERSON}/\texttt{VISIT\_OCCURRENCE}/\texttt{CONDITION\_OCCURRENCE}/
\texttt{MEASUREMENT} as features, as a first analytical/ML layer beyond
ETL and standards mapping; and extension to a wearable-native source (e.g.\
WESAD or PPG-DaLiA), mapping physiological signals into FHIR
\texttt{Observation}/\texttt{Device} and then into \texttt{MEASUREMENT}.

\section{Conclusion}

I presented a verified, open FHIR-to-OMOP CDM v5.4 ETL pipeline covering
four clinical resource types on the MIMIC-IV-FHIR demo dataset, with every
quantitative claim confirmed against a running database rather than
assumed. Beyond the implementation, I argue for decomposing
concept-mapping coverage metrics to separate source-data ceilings from
mapping-logic defects, and I describe a synthetic-fixture CI methodology
that tests real pipeline code under the real licensing constraints that
apply to any OMOP ETL project built on OHDSI's Athena vocabularies. Two
edge cases surfaced by deliberately adversarial test design, absent from
the real dataset only by chance, illustrate a broader point: an ETL
pipeline validated only against the data it happens to have seen carries
untested assumptions that a sufficiently different cohort could violate.

\section*{Availability}
\label{sec:availability}
Source code, SQL, tests, and CI configuration are publicly available at
\url{https://github.com/Saileshwar1625/fhir-omop-etl} (tag \texttt{v1.0.1}).

\bibliographystyle{plain}
\begin{thebibliography}{9}

\bibitem{hl7fhir}
HL7 International.
\emph{FHIR Release 4 (R4)}.
\url{https://hl7.org/fhir/R4/}.

\bibitem{hripcsak2015ohdsi}
Hripcsak G, Duke JD, Shah NH, Reich CG, Huser V, Schuemie MJ, Suchard MA,
Park RW, Wong ICK, Rijnbeek PR, van der Lei J, Pratt N, Norén GN, Li Y-C,
Stang PE, Madigan D, Ryan PB.
Observational Health Data Sciences and Informatics (OHDSI): Opportunities
for Observational Researchers.
\emph{Stud Health Technol Inform}. 2015;216:574--578.

\bibitem{bennett2023mimicfhir}
Bennett AM, Ulrich H, van Damme P, Wiedekopf J, Johnson AE.
MIMIC-IV on FHIR: converting a decade of in-patient data into an
exchangeable, interoperable format.
\emph{J Am Med Inform Assoc}. 2023;30(4):718--725.
\url{https://doi.org/10.1093/jamia/ocad002}.

\bibitem{bennett2025physionet}
Bennett A, Ulrich H, Wiedekopf J, Szul P, Grimes J, Johnson A.
MIMIC-IV Clinical Database Demo on FHIR (version 2.1.0).
\emph{PhysioNet}. 2025.
\url{https://doi.org/10.13026/vphg-y548}.

\bibitem{johnson2023mimiciv}
Johnson AEW, Bulgarelli L, Shen L, Gayles A, Shammout A, Horng S,
Pollard TJ, Hao S, Moody B, Gow B, Lehman L-wH, Celi LA, Mark RG.
MIMIC-IV, a freely accessible electronic health record dataset.
\emph{Sci Data}. 2023;10:1.
\url{https://doi.org/10.1038/s41597-022-01899-x}.

\bibitem{goldberger2000physionet}
Goldberger AL, Amaral LAN, Glass L, Hausdorff JM, Ivanov PCh, Mark RG,
Mietus JE, Moody GB, Peng C-K, Stanley HE.
PhysioBank, PhysioToolkit, and PhysioNet: Components of a New Research
Resource for Complex Physiologic Signals.
\emph{Circulation}. 2000;101(23):e215--e220.

\bibitem{ohdsi_cdm}
OHDSI.
\emph{OMOP Common Data Model}, version 5.4.0.
\url{https://github.com/OHDSI/CommonDataModel}.

\end{thebibliography}

\end{document}
